
\subsubsection{Interpretation}

Batch Normalization computes statistics over the batch dimension. When training batches inherit spurious correlations from the dataset (e.g., 90\% of landbird samples have grass backgrounds), Batch Normalization's batch-level mean and variance encode this spurious pattern. This amplifies gradients toward majority-group samples that align with the spurious correlation, accelerating shortcut learning.

Layer Normalization normalizes each instance independently over feature dimensions. It does not aggregate information across the batch dimension and therefore does not amplify batch-level spurious patterns.

The gradient measurement (RQ2) provides mechanistic evidence: Batch Normalization shows 26\% higher gradient flow to spurious-aligned samples during epochs 0--19, explaining why it reaches high average accuracy (majority groups dominate) while maintaining large worst-group gaps (minority groups receive weaker updates).

\subsubsection{Comparison with Group DRO}

Sagawa et al. (2020) improved Waterbirds worst-group accuracy from 72.6\% (ERM with ResNet-50-BN) to 91.4\% (Group DRO with ResNet-50-BN) by reweighting loss toward minority groups. Our approach reduces worst-group gaps via architectural choice (Layer Normalization vs Batch Normalization) under standard ERM loss, without requiring group labels. Direct comparison is not possible because our experiments used synthetic data, not real Waterbirds. The 9.41 pp gap reduction suggests that architectural choice and algorithmic interventions are complementary.

\subsubsection{Limitations}

\#\#\#\# L1: Synthetic Data

All results are based on synthetic spurious correlation data (90\% co-occurrence) due to WILDS Waterbirds server unavailability (HTTP 500 error). Effect sizes (d = 3.94) may be inflated compared to real datasets. Real Waterbirds validation is required before claiming generalization.

Impact: Results demonstrate workflow and hypothesis testing procedures. Quantitative claims (9.41 pp gap) may change by 20--50\% on real data.

\#\#\#\# L2: Attention Hypothesis Incomplete

The planned attention mechanism hypothesis (h-m2) was not completed due to: training process failure (1/6 runs), cuDNN initialization error, and dataset download failure. Conclusions are restricted to Batch Normalization vs Layer Normalization. Claims about attention mechanisms (CBAM, ViT) are deferred to future work.

Impact: Scope narrowed from ''normalization + attention'' to ''normalization only.''

\#\#\#\# L3: Vision Tasks Only

All experiments target vision tasks. Generalization to NLP, audio, or tabular data is unknown. Batch Normalization is less common in NLP (Transformers use Layer Normalization), suggesting the finding may be vision-specific.

Impact: Claims restricted to computer vision domain.

\#\#\#\# L4: Constant Learning Rate

Experiments used constant learning rate (0.01) to isolate architectural effects. Real-world training uses learning rate schedules (warmup, cosine decay). Whether the BN-LN gap persists under scheduled learning rates is unknown.

Impact: Results valid for constant learning rate regime (common in Group DRO literature). Generalization to scheduled learning rates untested.

\subsubsection{Broader Implications}

This study introduces accuracy-matched temporal comparison as a methodology for studying architectural effects on robustness. By comparing architectures at matched average accuracy (e.g., 90\%) rather than fixed epochs, training speed confounds are eliminated.

The finding suggests that normalization layer choice affects fairness metrics on datasets with spurious correlations. When deploying models on data with potential spurious correlations, Layer Normalization may reduce worst-group disparities compared to Batch Normalization, pending real dataset validation.
